\cotefile{91}
\scanpage{129}
\mspage{- 135-}
% archive sheet number: 128
\hfill n\textsuperscript{o} 316

a) L'algèbre $\overset{\handwritten{\sim}}{U}$ et l'application $\handwritten{\unsure{\tilde{\varrho}|\tilde{\rho}}}$ sont solutions d'un problème d'application universelle : pour tout $p$-homomorphisme $f$ de $g$ dans une algèbre associative $B$ sur $K$ (considérée comme $p$-algèbre de Lie), il existe un $K$-homomorphisme et un seul $f'$ de $\overset{\handwritten{\sim}}{U}$ dans $B$ (pour les structures d'algèbres associatives) tel que $f=f'\unsure{o|\circ}\handwritten{\tilde{\varrho}}$.

b) Montrer que si $(b_{\handwritten{\lambda}})_{\handwritten{\lambda}\handwritten{\in}\handwritten{\Lambda}}$ est une base de $g$ ($\handwritten{\Lambda}$ étant totalement ordonné), $\handwritten{\tilde{\varrho}}$ est injective, et que, si on identifie $x$ et $\handwritten{\tilde{\varrho}}(x)$ pour $x\handwritten{\in} g$, les éléments $\handwritten{\prod_{\lambda}}b_{\handwritten{\lambda}}^{\handwritten{\nu_{\lambda}}}$ (où les $\handwritten{\lambda}$ vont en croissant, où les $\handwritten{\nu_{\lambda}}$ sont nuls sauf un nombre fini d'entre eux, et où $0\handwritten{\leqslant\nu_{\lambda}}\handwritten{<}p$ pour tout $\handwritten{\lambda}$) forment une base de $\overset{\handwritten{\sim}}{U}$. (Identifiant canoniquement $g$ à un sous-module de $U$, on pose $\phi(x)=x^{p}-x^{[p]}$ pour tout $x\handwritten{\in} g$. Pour tout indice composé $\alpha=(\alpha_{\handwritten{\lambda}})\handwritten{\in}\unsure{\mathbb{N}|N}^{(L)}$, soit $\alpha_{\handwritten{\lambda}}=\handwritten{\beta_{\lambda}}+p\handwritten{\gamma_{\lambda}}$ avec $0\handwritten{\leqslant\beta_{\lambda}}\handwritten{<}p$, et soit $T_{\alpha}=(\handwritten{\prod_{\lambda}}b_{\handwritten{\lambda}}^{\handwritten{\beta_{\lambda}}})(\handwritten{\prod_{\lambda}}(\phi(b_{\handwritten{\lambda}}))^{\handwritten{\gamma_{\lambda}}})$. Montrer que les $T_{\alpha}$ forment une base de $U$ et les $T_{\alpha}$ tels que $\handwritten{\gamma}=(\handwritten{\gamma_{\lambda}})\ne 0$ une base de $\unsure{J|\mathscr{J}}$ ; on observera que les $\phi(b_{\handwritten{\lambda}})$ appartiennent au centre de $U$).

7) - Soit $g$ une $p$-algèbre de Lie sur un anneau $K$ de caractéristique $p>0$ ($p$ premier). On dit qu'une dérivation $D$ de $g$ est une $p$\uline{-dérivation} si on a $D(x^{p})=(\operatorname{ad}x)^{p-1}.dx$ pour tout $x\handwritten{\in} g$. Toute dérivation intérieure est une $p$-dérivation.

a) Si $L$ est une $K$-algèbre associative, toute dérivation de $L$ est une $p$-dérivation lorsque $L$ est considérée comme $p$-algèbre de Lie. (Utiliser la formule (2) de l'exerc. 19 du \S 1).

b) Supposons que $g$ ait une base. Pour qu'une dérivation de $g$ soit une $p$-dérivation, il faut et il suffit qu'elle puisse se prolonger en une dérivation de l'algèbre enveloppante restreinte de $g$.
